\begin{explanationbox}
	\textbf{\textcolor{CExplanation}{一句话直觉}}
	万能公式通过引入半角正切 $t = \tan\frac{x}{2}$，把 $\sin x$、$\cos x$、$\tan x$ 统一表示为 $t$ 的有理函数，从而将三角问题转化为代数问题。

	\medskip
	\textbf{\textcolor{CExplanation}{核心拆解}}
	\begin{description}[style=nextline, leftmargin=2.8em, labelsep=0.5em, itemsep=0.45em]
		\item[\textbf{要点一（半角代换）：}]
			令 $t = \tan\frac{x}{2}$，把原变量 $x$ 替换为 $t$。
		\item[\textbf{要点二（二倍角转化）：}]
			利用 $\sin x = 2\sin\frac{x}{2}\cos\frac{x}{2}$，$\cos x = \cos^2\frac{x}{2} - \sin^2\frac{x}{2}$。
		\item[\textbf{要点三（有理化处理）：}]
			将分子分母同除以 $\cos^2\frac{x}{2}$，转化为关于 $\tan\frac{x}{2}$ 的有理式。
	\end{description}

	\medskip
	\textbf{\textcolor{CExplanation}{几何本质（半角投影）}}
	考虑单位圆上点 $P(\cos x, \sin x)$ 与定点 $A(-1,0)$ 的连线，其斜率恰好等于 $\tan\frac{x}{2}$，由此可导出 $\sin x$、$\cos x$ 与 $t$ 的上述关系。

	\medskip
	\textbf{\textcolor{CExplanation}{代数意义（有理化简）}}
	万能公式将三角函数的复合运算（和差化积、倍角等）转化为有理式的加减乘除，降低技巧门槛，适合程序化求解。

	\medskip
	\textbf{\textcolor{CExplanation}{考点价值（应用场景）}}
	\begin{itemize}[leftmargin=*, itemsep=0.5em, topsep=0.4em]
		\item \textbf{考法一（求值计算）：} 已知 $\tan\frac{x}{2}$ 时，直接代入公式求 $\sin x,\cos x,\tan x$。
		\item \textbf{考法二（恒等证明）：} 将待证式化为 $t$ 的有理式，证明等价性。
		\item \textbf{考法三（方程求解）：} 将三角方程转化为关于 $t$ 的代数方程，解出 $t$ 后再反求 $x$。
	\end{itemize}

	\medskip
	\textbf{\textcolor{CExplanation}{顿悟点}}
	\textbf{万能公式分母总是 $1+t^2$，恒为正，因此不会出现分母为零导致无意义的情况（正切公式除外）。}

	\medskip
	\textbf{\textcolor{CExplanation}{使用场景}}
	\begin{itemize}[leftmargin=*, itemsep=0.5em, topsep=0.4em]
		\item \textbf{场景一（已知半角正切）：} 题目直接给出 $\tan\frac{x}{2}$ 的值。
		\item \textbf{场景二（分式化简）：} 含 $\sin x,\cos x$ 的齐次分式可化为 $t$ 的有理式。
		\item \textbf{场景三（非特殊角求解）：} 需要计算某一三角函数值但无直接公式时，可设 $t$ 转化。
	\end{itemize}

\end{explanationbox}
